\chapter{Advanced Time Series and Panel Dynamics}

\section{Time-varying parameters}

In \texttt{tvp} the coefficient vector is a random walk observed with noise:

\[
y_t = x_t'\beta_t + \varepsilon_t, \qquad
\beta_t = \beta_{t-1} + \eta_t
\]

\texttt{tvp} estimates the variances of $\varepsilon_t$ and $\eta_t$ by MLE and
reports the smoothed path of $\beta_t$.

\begin{lstlisting}[caption={Time-varying parameter regression}]
let n = 200
generate df x = rnormal(n)
generate df y = (1 + 0.02*_n) + (0.5 - 0.003*_n)*x + rnormal(n)

let m_tvp = tvp(y ~ x, df)
print(m_tvp)
\end{lstlisting}

Use \texttt{tvp\_var} for a full vector autoregression with time-varying
coefficients.

\section{Threshold and Markov-switching VAR}

\texttt{setar} allows two AR regimes split by a threshold value of the lagged
dependent variable.

\begin{lstlisting}[caption={SETAR}]
let m_setar = setar(y, df, order=2, delay=1)
print(m_setar)
\end{lstlisting}

\texttt{msvar} allows the coefficients themselves to switch according to a
latent Markov chain.

\begin{lstlisting}[caption={Markov-switching VAR}]
let m_msvar = msvar(y1 + y2, df, regimes=2, lags=1)
print(m_msvar)
\end{lstlisting}

\section{Asymmetric and nonlinear dynamics}

\texttt{nardl} decomposes a regressor into positive and negative partial sums,
allowing asymmetric long-run responses.

\begin{lstlisting}[caption={NARDL}]
let m_nardl = nardl(y ~ x, df, lags=1)
print(m_nardl)
\end{lstlisting}

\texttt{pstr} does the same in a panel, with coefficients that change smoothly
according to a logistic transition function of an observed variable.

\begin{lstlisting}[caption={Panel Smooth Transition Regression}]
let m_pstr = pstr(y ~ x1 + x2, df, q="z", id="firm")
print(m_pstr)
\end{lstlisting}

\section{Functional coefficients}

\texttt{fcoef} estimates coefficients that vary smoothly with a moderator
variable $z$, using local linear kernel regression.

\begin{lstlisting}[caption={Functional coefficient model}]
let m_fcoef = fcoef(y ~ x1 + x2, df, z="z", points=20)
print(m_fcoef)
\end{lstlisting}

\section{Bayesian VAR and mixed-frequency VAR}

\texttt{bvar} estimates a VAR with a Minnesota prior, shrinking own lags
towards 1 and cross-lags towards 0.

\begin{lstlisting}[caption={Bayesian VAR}]
let m_bvar = bvar(y1 ~ y2, df, lags=1)
print(m_bvar)
\end{lstlisting}

\texttt{mfvar} combines low- and high-frequency variables using an exponential
Almon MIDAS aggregation.

\begin{lstlisting}[caption={Mixed-frequency VAR}]
let m_mfvar = mfvar(df_q, yq, df_m, ym, agg=3, lags=1)
print(m_mfvar)
\end{lstlisting}

\texttt{fcoef}, \texttt{bvar} and \texttt{mfvar} are experimental and
require \texttt{--features experimental}.

\section{Panel VAR}

\texttt{pvar} estimates a panel VAR using Arellano-Bond style GMM: lagged
levels instrument first-differenced equations.

\begin{lstlisting}[caption={Panel VAR}]
let m_pvar = pvar(y1 ~ y2, df, id="firm", lags=1)
print(m_pvar)
\end{lstlisting}

These commands are powerful but convention-sensitive. Always compare the
output with a reference implementation when possible and inspect the
transition probabilities, eigenvalues, or J-statistics reported by the
estimator.
